\documentclass[12pt]{article}
\usepackage[utf8]{inputenc}
\usepackage[T1]{fontenc}
\usepackage{hyperref}
\title{Zenag Language Resources for K-12 Curriculum Development}
\author{Compiled by Claude Sonnet 5.5 (claude-sonnet-5-5) for Andr\'as Kornai}
\date{}
\begin{document}
\maketitle

\section{Introduction}
Zenag (ISO 639-3 zeg) is an Austronesian language of the Huon Gulf branch, spoken in the Mumeng area of Morobe Province, Papua New Guinea. Speaker estimates vary widely by source, from about 1,900 (a 1988 SIL figure) to about 7,000. Speakers also use Tok Pisin. Very little documentation turned up in the searches.

\section{Primary Educational Resources}
No primers, textbooks or school programmes were found.

\section{Linguistic and Reference Materials}
\begin{itemize}
\item Mumeng language, Wikipedia: Zenag is treated as part of this dialect chain, with a 1979 combined figure of 7,100 speakers: \url{https://en.wikipedia.org/wiki/Mumeng_language}
\item Numeral-systems page on Zenag, Max Planck Institute for Evolutionary Anthropology: \url{https://lingweb.eva.mpg.de/channumerals/Zenag.htm}
\item Endangered Languages Project entry, based on Ethnologue 16th edition (2009): \url{https://www.endangeredlanguages.com/lang/2704?hl=en}
\end{itemize}

\section{Religious and Community Resources}
Joshua Project reports no Scripture translation and no written form recorded for Zenag: \url{https://joshuaproject.net/people_groups/16098/PP}

\section{Digital and Audio Resources}
No audio, video or corpus resources were found.

\section{Curriculum Development Context}
Zenag has no known written materials, so a curriculum would start from new documentation and orthography work in collaboration with the community. Tok Pisin is a likely bridging language. The Joshua Project language page gives basic population data: \url{https://joshuaproject.net/index.php/languages/zeg}
\end{document}